\section{Método de verificação}
\label{sec:metodo}

\status{Pronta. O catálogo de soluções, as normas e o oráculo de resíduo do
operador estão todos implementados e são o que produziu os números da seção
\ref{sec:resultados}.}

\subsection{O problema e as soluções manufaturadas}

Resolve-se, no quadrado unitário $\Omega=[0,1]^2$,
\begin{equation}
  \nabla^2 u = f \quad\text{em } \Omega,
  \qquad
  \alpha u + \beta\,\frac{\partial u}{\partial n} = g \quad\text{em } \partial\Omega .
  \label{eq:poisson}
\end{equation}

O método das soluções manufaturadas inverte a direção usual: escolhe-se $u$,
\emph{deriva-se} $f=\nabla^2 u$ analiticamente, impõem-se as condições de
contorno a partir do próprio $u$, e compara-se a solução discreta com $u$. A
vantagem decisiva é que a solução exata é conhecida em todo ponto, e portanto o
erro é mensurável sem refino de referência.

O catálogo tem seis soluções, com todas as derivadas escritas a mão --- nunca por
diferenciação numérica, que introduziria o erro que se quer medir:

\begin{center}
\begin{tabular}{lll}
\toprule
nome & grau & papel \\
\midrule
\texttt{grau1} & 1 & campo linear: todo esquema consistente é exato \\
\texttt{grau2} & 2 & o pedido original; separa malha uniforme de adaptada \\
\texttt{grau3} & 3 & separa o fluxo de segunda ordem do compacto e do MLS \\
\texttt{grau4} & 4 & separa o compacto do MLS de ordem 4 \\
\texttt{grau6} & 6 & acima de toda exatidão disponível \\
\texttt{trig}  & --- & $\sin/\cos$: não polinomial, logo há ordem a medir \\
\bottomrule
\end{tabular}
\end{center}

Os coeficientes de cada monômio são \textbf{distintos}, de propósito: assim o
valor identifica a posição, e um esquema que troque os eixos não passa por acaso.

\subsection{As normas e a ordem}

O erro é medido em duas normas, a do máximo e a $L_2$ ponderada por área:
\begin{equation}
  \|e\|_\infty = \max_c |u_c - u(\mathbf{x}_c)|,
  \qquad
  \|e\|_2 = \sqrt{\frac{\sum_c A_c\,(u_c - u(\mathbf{x}_c))^2}{\sum_c A_c}} .
  \label{eq:normas}
\end{equation}
A ponderação por área da norma $L_2$ é necessária em malha adaptada: sem ela, as
células pequenas --- que são muitas --- dominariam a média sem corresponder à
fração do domínio que ocupam.

A ordem de convergência é estimada entre níveis consecutivos de refino,
\begin{equation}
  p = \log_2 \frac{\|e\|_{h}}{\|e\|_{h/2}} .
  \label{eq:ordem}
\end{equation}

\paragraph{Onde o esquema é exato não há ordem para medir.} Isso não é uma
ressalva: é a razão pela qual as tabelas da seção \ref{sec:resultados} exibem um
travessão em vez de um número. Com erro de máquina em dois níveis consecutivos,
\eqref{eq:ordem} divide ruído por ruído e produz qualquer valor. Imprimir esse
valor sugeriria uma taxa de convergência onde existe apenas aritmética de ponto
flutuante.

\subsection{O piso do erro é o solver}

Todas as corridas usam PETSc com tolerância \num{1e-14}. O erro residual
observado onde o esquema é exato fica na faixa de \num{1e-15} a \num{1e-13}, e
ele é \textbf{resíduo do solver linear}, não do esquema. Exigir menos que isso
seria pedir ao solver algo que ele não promete com essa tolerância. Os casos da
suíte usam \num{1e-10} como limiar de ``erro de máquina'' por essa razão.

\subsection{O resíduo do operador, como oráculo independente}
\label{sec:met-residuo}

A comparação da solução discreta com a exata mede o sistema inteiro: esquema,
condições de contorno e solver juntos. Para separar o esquema do resto, há um
segundo oráculo, que não resolve nada:
\begin{equation}
  r = \max_c \frac{\bigl|(A\,\mathbf{u}^{\text{exata}})_c - b_c\bigr|}{\text{normalização}_c},
  \label{eq:residuo}
\end{equation}
isto é: monta-se o operador, aplica-se à \emph{solução exata amostrada nos centros
das células}, e compara-se com o lado direito. Se o esquema é exato para aquela
solução, $r$ é de máquina --- sem solver envolvido.

Duas propriedades fazem desse oráculo mais que uma conveniência. Primeiro, ele
distingue \textbf{consistência} de \textbf{convergência}: um operador pode ser
consistente e o sistema divergir por causa do contorno, ou vice-versa. Segundo, a
montagem é \emph{a mesma função} que o solver usa --- extraída para que os dois
caminhos compartilhem um único código. Duas montagens separadas divergiriam
exatamente onde é difícil ver: perto da parede.

O valor desse oráculo ficou patente na transição 2:1. O resíduo do operador em
malha adaptada, para grau~2, dá \num{4.5} --- contra \num{2.27e-13} na uniforme.
Esse contraste de treze ordens de grandeza é o que afirma, sem ambiguidade, que a
diferença é do \emph{esquema na transição}, não do solver nem do contorno.

\subsection{Contraste obrigatório}

Nenhuma afirmação de exatidão entra na suíte sozinha. Para cada ``o erro é de
máquina'' há um caso vizinho medindo a condição em que ele deixa de ser:

\begin{center}
\begin{tabular}{ll}
\toprule
afirmação & contraste que a acompanha \\
\midrule
grau 2 exato com fechamento de 3 pontos & grau 2 com fechamento de 2 pontos \\
grau 3 exato no MLS de ordem 3 & grau 3 no fluxo de segunda ordem \\
grau 2 exato em malha adaptada com MLS & grau 2 em malha adaptada com fluxo \\
grau 4 exato com Gauss $2\times2$ na fonte & o mesmo com a fonte no ponto médio \\
pesos próprios concordam com o herdado & concordância \emph{exatamente} nula é recusada \\
\bottomrule
\end{tabular}
\end{center}

A última linha merece explicação: um caso que afirma ``as duas implementações
concordam'' e obtém diferença exatamente zero não está comparando duas
implementações --- está comparando uma consigo mesma. Esse foi um defeito real na
comparação CPU--GPU (seção \ref{sec:est-gpu}), e desde então a concordância
perfeita é tratada como sintoma, não como sucesso.
